\section{Arithmetic-step Gamma and Hurwitz--Stieltjes bridges}
\label{sec:hurwitz}

The rational interpolation of a coefficient column gives a direct link to Gamma products and ordinary Stieltjes primitives. This complements the power-series polylogarithmic lift.

\subsection{Step sums and Gamma products}

Let $\rho(x)=\sum_jd_j/(x+\alpha_j)$ be the rational function in \eqref{eq:rho}. For $h>0$ and real $x> -\min_j\alpha_j$, put
\[
 S_h(x)=\sum_{n\ge0}\rho(hn+x).
\]
All identities below also continue analytically away from their poles, with compatible logarithm branches.

\begin{proposition}[Gamma-product primitive]\label{prop:steps}
The sum converges absolutely and
\begin{align}
 S_h(x)&=-\frac1h\sum_jd_j\psi\left(\frac{x+\alpha_j}{h}\right),\label{eq:steps}\\
 \frac{\mathrm d^q}{\mathrm dx^q}S_h(x)
 &=-\frac1{h^{q+1}}\sum_jd_j\psi^{(q)}\left(\frac{x+\alpha_j}{h}\right),
 \qquad q\ge0,\label{eq:steps-derivative}\\
 \int_{x_0}^xS_h(t)\dd t
 &=-\sum_jd_j\log\frac{\Gamma((x+\alpha_j)/h)}
                              {\Gamma((x_0+\alpha_j)/h)}.\label{eq:steps-primitive}
\end{align}
\end{proposition}
\begin{proof}
The moment cancellations \eqref{eq:moment-cancel} give $\rho(x)=O(x^{-4})$ and justify the convergent sum. For a finite cutoff $N$, the digamma recurrence gives
\[
 \sum_{n=0}^{N}\frac1{hn+x+\alpha_j}
 =\frac1h\left\{\psi\left(N+1+\frac{x+\alpha_j}{h}\right)
                -\psi\left(\frac{x+\alpha_j}{h}\right)\right\}.
\]
The common logarithmic divergence cancels because $\sum d_j=0$, and the remaining upper-end difference tends to zero. This proves \eqref{eq:steps}. Differentiation of the normally convergent series, or of its finite digamma expression, proves \eqref{eq:steps-derivative}; integration gives \eqref{eq:steps-primitive} with its exact anchor.
\end{proof}

For $h=2$ and the fixed parity shift, the arguments differ by integers, so repeated use of the digamma recurrence reduces the sum to a rational function. Other arithmetic steps can retain genuine rational-argument Gamma and polygamma values.

For the stripped mixed weight $R$ from \eqref{eq:Rexample}, the same calculation gives
\begin{align}
 \sum_{n\ge0}R(3n)&=-\frac{11}{384}+\frac{\log3}{128}
                         +\frac{5\sqrt3\,\pi}{1152},\label{eq:Rstep3}\\
 \sum_{n\ge0}R(4n)&=\frac{13}{360}-\frac{\pi}{96}.\label{eq:Rstep4}
\end{align}
To verify these values, substitute the four coefficients of \eqref{eq:Rpartial} in \eqref{eq:steps}. Reduce shifts using $\psi(z+1)=\psi(z)+1/z$ and then apply the reflection and multiplication formulas at thirds or quarters. The accompanying exact script independently carries out these finite reductions. Equation~\eqref{eq:sumR} is the corresponding step-two case.

\subsection{A balanced entire Hurwitz family}

For the same fixed coefficients and $a_j(x)=(x+\alpha_j)/h$, define
\begin{equation}\label{eq:Zbalanced}
 Z(s,x)=\sum_jd_j\zeta(s,a_j(x)).
\end{equation}
The Hurwitz zeta function has only its usual simple pole at $s=1$, with residue one. Hence $\sum d_j=0$ makes $Z$ entire in $s$. The stronger cancellations in \eqref{eq:moment-cancel} also give
\begin{equation}\label{eq:Zzeros}
 Z(0,x)=Z(-1,x)=0.
\end{equation}
Indeed the special values at $0$ and $-1$ are polynomials of degrees one and two in $a_j(x)$, and the three moment cancellations annihilate them. These are exact identities, not assertions about zeros of an individual Hurwitz zeta function.

Use the Stieltjes convention
\begin{equation}\label{eq:stieltjes}
 \zeta(1+s,a)=\frac1s+\sum_{q\ge0}\frac{(-1)^q}{q!}\gamma_q(a)s^q.
\end{equation}
The balanced coefficients
\[
 B_q(x)=\sum_jd_j\gamma_q(a_j(x))
\]
therefore satisfy
\begin{equation}\label{eq:balanced-Laurent}
 Z(1+s,x)=\sum_{q\ge0}\frac{(-1)^q}{q!}B_q(x)s^q,
 \qquad B_0(x)=hS_h(x).
\end{equation}
The last equality uses $\gamma_0(a)=-\psi(a)$.

\begin{proposition}[All-index ordinary Stieltjes primitives]\label{prop:balanced-primitives}
For every $q\ge0$,
\begin{equation}\label{eq:balanced-primitive}
 \int_{x_0}^x B_q(t)\dd t
 =\frac{(-1)^{q+1}h}{q+1}
       \left[\partial_s^{q+1}Z(s,x)\big|_{s=0}
             -\partial_s^{q+1}Z(s,x_0)\big|_{s=0}\right].
\end{equation}
In particular
\begin{equation}\label{eq:ZprimeGamma}
 \partial_sZ(0,x)=\sum_jd_j\log\Gamma(a_j(x)),
\end{equation}
so the case $q=0$ is exactly the Gamma-product primitive in Proposition~\ref{prop:steps}.
\end{proposition}
\begin{proof}
The parameter derivative identity $\partial_a\zeta(s,a)=-s\zeta(s+1,a)$ gives
\[
 \partial_xZ(s,x)=-\frac{s}{h}Z(s+1,x).
\]
Insert \eqref{eq:balanced-Laurent} and compare the coefficient of $s^{q+1}$. This yields
\[
 \partial_x\partial_s^{q+1}Z(0,x)
 =-\frac{(-1)^q(q+1)}hB_q(x),
\]
which integrates to \eqref{eq:balanced-primitive}. Finally,
$\zeta'(0,a)=\log\Gamma(a)-\frac12\log(2\pi)$ and $\sum d_j=0$ give \eqref{eq:ZprimeGamma}.
\end{proof}

These Hurwitz parameter identities are classical \cite{DLMFHurwitz}. Here they provide an explicitly normalized bridge from the new universal coefficient densities to Gamma products and every Stieltjes order. Neither a finite part in the spatial variable nor a hidden additive constant occurs in \eqref{eq:balanced-primitive}.
